\paragraph{A locally split cup lift at auxiliary primes.}
\label{par:count-two-cup}
To join the two endpoints we next construct the required extension classes. The
residual path begins over the tame real field \(F_0\).
Write \(\Gamma=G_{F_0}\), let \(k\) be the common residual
coefficient field, and put
\[
 A=d_1/d_0,\qquad B=d_2/d_0 ,
\]
where \(d_i\) are the four ordered residual characters. Their
symplectic products satisfy \(d_0d_3=d_1d_2\). The normalizing
element \(\tau\) has distinct diagonal entries.

In particular
\(A(\tau)\ne1\) and \((AB)(\tau)\ne1\). The normalized target
has support in the diagonal and the entries
\[
 (1,0),\ (2,0),\ (3,0),\ (3,1),\ (3,2),\ (1,2).
\]
Its local lower unipotent symplectic frame \(N_v\) has
\((N_v)_{21}=0\). Put \(y(g)=r(g)_{20}/d_0(g)\).
Multiplication in that support gives
\[
 y(gh)=y(g)+B(g)y(h),\qquad y(\tau)=0 .
\]
For \(s\) in the marked decomposition group \(G_v\), the first
column of \(r(s)N_v=N_v(N_v^{-1}r(s)N_v)\) gives
\[
                    y(s)=(1-B(s))n_v,\qquad n_v=(N_v)_{20}.
\]
This uses the prescribed first diagonal character in the local
upper triangular matrix. It applies even when \(y\) is zero and
the local entry \(n_v\) is nonzero.

\begin{lemma}[The two rational shapes after restriction]\label{lem:count-two-target-restriction}
In the first lattice shape of Proposition~\ref{prop:rational-rank-four-torus},
the two entry functions on the restricted group satisfy
\[
 r_{20}|_{G_{F_0}}\not\equiv0,\qquad r_{12}|_{G_{F_0}}\not\equiv0 .
\]
In the second shape the global \((2,0)\) entry is zero, the normalized
marked frame has nonzero \((2,0)\) coordinate, and a \((1,0)\) entry is
nonzero on \(G_{F_0}\).
\end{lemma}

\begin{proof}
The fitted field \(F_0\) is abelian over \(\Q\). For an entry
whose normalized coordinate is a cocycle \(f\) of
character \(\theta\), suppose \(f\) vanishes on \(G_{F_0}\).
For any \(g\in G_\Q\), normality puts
\(g\tau g^{-1}\) in \(G_{F_0}\), and the cocycle
identity gives
\[
              \begin{aligned}
 0=f(g\tau g^{-1}) &= f(g)+\theta(g)f(\tau)
                       +\theta(g)\theta(\tau)f(g^{-1})\\
                   &=(1-\theta(\tau))f(g),
 \end{aligned}
\]
Since \(1-\theta(\tau)\ne0\), the restriction of a
nonzero normalized cocycle to \(G_{F_0}\) is nonzero.
The \((2,0)\) coordinate has character \(d_2/d_0\). For the
\((1,2)\) coordinate, no intermediate product contributes to
the \((1,2)\) entry, so dividing by \(d_2\) gives a cocycle
of character \(d_1/d_2\). The two coordinates therefore survive
in the first shape. The other allowed pattern is distinguished by the global representation entry and the local frame entry
\[
 r(g)_{20}=0\ (\text{all }g),\qquad n=N_{20}\ne0 .
\]

In this support the normalized \((1,0)\) entry is a
cocycle, since its only additional product would use
the zero \((2,0)\) entry. If it vanishes on the tower
it vanishes over \(\Q\) by the displayed argument.
On the rational marked decomposition group the
first-column equations in rows two and one then
give \(d_2=d_0\) and
\[
             r(s)_{12}=\frac{N_{10}}{n}(d_2(s)-d_1(s)).
\]
This contradicts the non-coboundary in the second
lattice shape. Thus the restricted point has a
nonzero \(-w_\alpha\) entry. The frame transport calculation in
Paragraph~\ref{par:count-two-restriction} gives
\(n\mapsto (d_2(g)/d_0(g))n\) when the global \((2,0)\) entry
vanishes. Thus the other coordinate required in the second
shape survives as well.
\end{proof}

\begin{lemma}[A prescribed tame ramification direction]\label{lem:cup-ramification-direction}
Let \(\ell\) be an odd prime, \(F_0\) a number field, and \(M\) a
finite continuous \(\mathbb F_\ell\)-representation of \(G_{F_0}\).
Fix a nonzero \(e_0\in M\), a finite set \(S\) containing the places
above \(\ell\), the infinite places and the ramification of \(M\),
and a finite Galois extension of \(\Q\) containing the normal closure
of \(F_0\), a splitting field of \(M\), and \(\mu_\ell\). There is a
continuous cocycle \(c\colon G_{F_0}\to M\) unramified outside \(S\)
and the places above finitely many rational primes splitting completely
in that extension. Its local class is zero at every \(v\in S\).
On the new inertia groups its image satisfies 
\[
 c(I_v)\subseteq\mathbb F_\ell e_0\quad(v\text{ new}),\qquad c(I_v)\ne0\quad\text{for at least one new }v.
\]

\end{lemma}

\begin{proof}
\proofstep{Specify the new local ramification.}
Take a set \(P\) of places over sufficiently many distinct
new rational primes, including all the places over each of those
primes. Chebotarev permits arbitrarily large \(|P|\).
Let \(M'=\operatorname{Hom}(M,\mu_\ell)\) with its natural Galois
action. At \(v\in P\) both modules are trivial and the residue
cardinality is one modulo \(\ell\). The quotient of
\(H^1(F_{0,v},M)\) by its unramified subspace identifies with the
homomorphisms from the cyclic exponent-\(\ell\) quotient of tame
inertia to \(M\). This follows from the tame presentation
\(\operatorname{Frob}\,t\,\operatorname{Frob}^{-1}=t^{q_v}\).
Requiring the image to lie in \(\mathbb F_\ell e_0\) gives a
one-dimensional subspace of this quotient.
\proofstep{Bound the dual obstruction.}
In the sum of local groups over \(S\cup P\) take the subgroup whose component at each place is
\[
 U_v=\begin{cases}H^1_{\rm unr}(G_v,M)&v\in P,\\0&v\in S.\end{cases}
\]
By local Tate duality its annihilator imposes unramifiedness at
each place of \(P\). By Poitou--Tate exactness the annihilator of
the image of \(H^1(G_{F_0,S\cup P},M)\) is the image of
\(H^1(G_{F_0,S\cup P},M')\). Intersect this latter image with
\(U^\perp\). Any representative global class is unramified at
\(P\), since the local inertia actions are trivial there and
coboundaries vanish on those inertia groups.

It inflates from
\(H^1(G_{F_0,S},M')\), by factoring the cocycle through the quotient
that kills the normal closure of the new inertia groups.
This fixed cohomology group is finite. Let \(L\) be the
localization image modulo \(U\), in the quotient of the sum
\(V_{\rm loc}\) of local groups over \(S\cup P\). The annihilator
description gives
\[
 \operatorname{codim}_{V_{\rm loc}/U}L
 =\dim L^\perp\le \dim_{\mathbb F_\ell}H^1(G_{F_0,S},M').
\]
\proofstep{Obtain a class in the prescribed directions.}
Let \(D_P\subseteq V_{\rm loc}/U\) be the direct sum of the
prescribed ramification lines and zero entries at \(S\).
Then \(\dim D_P=|P|\). Choose
\(|P|>\dim_{\mathbb F_\ell}H^1(G_{F_0,S},M')\).
The dimension formula gives
\[
 \dim(D_P\cap L)\ge |P|-\operatorname{codim}_{V_{\rm loc}/U}L>0.
\]
Lift a nonzero element of this intersection to a global class.
At the new places the inertia restriction is independent of its
representative and is nonzero on some inertia element, which acts
trivially on \(M\). At an old place, a representative is a
coboundary on the full decomposition group, so is zero on any
subgroup on which \(M\) acts trivially.
\hfill\mbox{\cite{milne-duality}}\label{cor:cohomology-count-two-chain-1}
\end{proof}

\begin{proposition}[The marked auxiliary symplectic datum]\label{prop:count-two-auxiliary-datum}
Take the residual characters and the marked \((2,0)\) target
column obtained by restriction in
Lemma~\ref{lem:count-two-target-restriction}. For these data,
one can choose an auxiliary representation \(r'\) with the same diagonal.
Writing \(r,N_v\) for the target matrices and frames and \(N'_v\) for the auxiliary frames, one has
\[
 r'_{20}=r_{20},\qquad (N'_v)_{20}=(N_v)_{20}\quad(v\text{ marked}),
                         \qquad r'_{10}\not\equiv0 .
\] It is diagonal at
\(\tau\), kills the required old inertia and is locally conjugate
to the marked diagonal at the prescribed local places. The new
ramification is supported above finitely many completely split
rational primes as in Lemma~\ref{lem:cup-ramification-direction}.
\end{proposition}

\begin{proof}
\proofstep{Lift the new root together with its cup correction.}
Apply Lemma~\ref{lem:cup-ramification-direction} with \(M\)
carrying coordinates \((x,z)\)
and action
\[
 \rho(g)(u,v)=\bigl(A(g)u,\ (AB)(g)v+y(g)A(g)u\bigr).
\]
The second component of \(\rho(g)\rho(h)(u,v)\) is
\[
 (AB)(gh)v+
 A(gh)\bigl(B(g)y(h)+y(g)\bigr)u .
\]
This is the second component of \(\rho(gh)(u,v)\) by the cocycle
identity for \(y\). Let \(k_{\rm d}\) be the finite subfield generated
by the diagonal characters' images. The span \(V_y\) over
\(k_{\rm d}\) of the finite image of \(y\) is finite, since \(y\) is
continuous on a compact group with discrete target. The chosen additive module is finite even when \(k\) is infinite, since it is the sum of the two finite pieces in the action:
\[
 M=k_{\rm d}\oplus V_y,\qquad |M|=|k_{\rm d}|\,|V_y| .
\]

Take \(e_0=(1,0)\).
All originally prescribed ramification is retained by taking
\(S\) large enough to contain it and using the same construction
with the additional condition that the rational primes split in
the field cut out by the residual diagonal, the quadratic CM
character and the mod-\(\ell\) cyclotomic character. On any old
inertia on which the former representation is required to be
the identity both the diagonal and \(y\) are trivial. The local
coboundary constructed there is zero on that subgroup.

The cocycle coordinates obey
\[
\begin{aligned}
 x(gh)&=x(g)+A(g)x(h),\
 z(gh)&=z(g)+(AB)(g)z(h)+y(g)A(g)x(h).
\end{aligned}
\]
On \(G_v\) write the local coboundary as
\[
 x(s)=(1-A(s))m_v,\qquad
 z(s)=(1-(AB)(s))o_v-y(s)A(s)m_v .
\]
Subtract a global coboundary to make \(x(\tau)=z(\tau)=0\).
At \(\tau\) the action is diagonal because \(y(\tau)=0\), and
neither of its two eigenvalues is one. Thus the subtraction is
possible and just changes the local constants. At the inertia element where \(x\ne0\), the action is trivial, so the subtraction leaves \(x\) unchanged. The same argument works with finitely
many fixed \(y_j\), by replacing \(z\) by a tuple with the indicated
action in each coordinate.

\proofstep{Produce the symplectic matrices and local flags.}
Set
\[
\begin{aligned}
 U(x,y,t)&=
 \begin{pmatrix}
 1&0&0&0\\ x&1&0&0\\ y&0&1&0\\ t&y&-x&1
 \end{pmatrix},\\
 U(x,y,t)U(x',y',t')&=U(x+x',y+y',t+t'+yx'-xy') .
\end{aligned}
\]
For the alternating form pairing positions \(0,3\) and \(1,2\)
with the same sign, \(U\) is symplectic and conjugation by
\(\mathcal D(g)=\operatorname{diag}(d_i(g))\) scales its three
coordinates by \(A(g),B(g),(AB)(g)\).
Let
\[
 R_1(g)=U(x(g),y(g),2z(g)-x(g)y(g))\mathcal D(g).
\]
Writing \(c=2z-xy\), the two cocycle equations give
\[
 c(gh)-c(g)-(AB)(g)c(h)
   =y(g)A(g)x(h)-x(g)B(g)y(h).
\]
Together with the identities for \(x,y\) this is exactly the
multiplication law for the unipotent factor of
\(R_1(g)R_1(h)\). Thus \(R_1\) is a representation with
the prescribed similitude and diagonal.

It is diagonal at \(\tau\), kills the required inertia, and has a
nonzero \((1,0)\) entry. At \(v\) take
\[
                         N_{1,v}=U(m_v,n_v,2o_v-m_vn_v).
\]
On \(G_v\), conjugating \(\mathcal D(s)\) by this matrix gives
\(R_1(s)\). In the third unipotent coordinate both sides give
\((1-(AB)(s))(2o_v-m_vn_v)+(B(s)-A(s))m_vn_v\).
This proves the flag conditions with both the full local
diagonal and the smaller inertia quotient. It also retains
the given \((2,0)\) matrix and frame entries.
\label{cor:cohomology-count-two-chain-2}
\end{proof}

\paragraph{Removing the auxiliary ramification and transporting the flags.}
\label{par:count-two-restriction}
The cup lift introduced tame ramification at the new primes. Before comparing the flags we restrict to a compositum of real cyclotomic subfields on which the new inertia cocycles vanish. The root classes needed for the path must survive this same restriction.

\begin{auxlemma}[Restriction of a normalized root cocycle]
\label{lem:root-cocycle-ell-power}
Let \(k\) be a field of characteristic \(\ell\), \(G\) a group,
\(\theta:G\to k^\times\) a character and \(u:G\to k\) a nonzero
cocycle for the action by \(\theta\). Suppose \(u(\tau)=0\) and
\(\theta(\tau)\ne1\) for some \(\tau\in G\). For any normal subgroup \(H\) the nonvanishing test is 
\[
 \exp(G/H)\mid\ell^a\quad\Longrightarrow\quad u|_H\ne0.
\]

\end{auxlemma}

\begin{proof}
Choose \(g\) with \(u(g)\ne0\). When \(\theta(g)=1\), replace \(g\)
by \(\tau g\). Its cocycle value is \(\theta(\tau)u(g)\ne0\) and
its character value is \(\theta(\tau)\ne1\).
For the resulting \(g\) the cocycle identity and the Frobenius map give
\[
 u(g^{\ell^a})=u(g)\sum_{j=0}^{\ell^a-1}\theta(g)^j
 =u(g)\frac{\theta(g)^{\ell^a}-1}{\theta(g)-1}
 =u(g)(\theta(g)-1)^{\ell^a-1}\ne0 .
\]
The element \(g^{\ell^a}\) belongs to \(H\).
\end{proof}

\begin{proposition}[Restriction and transport of the marked flags]\label{prop:count-two-tame-return}
For each new rational prime \(w\) of the auxiliary datum of
Proposition~\ref{prop:count-two-auxiliary-datum}, adjoin the real
cyclic degree-\(\ell\) subfield of \(\mathbb Q(\mu_w)\) to \(F_0\).
Let \(F_1\) be the resulting compositum. For each of these new
places \(v\) of \(F_0\), with tame inertia cocycle \(z_v\), restriction gives
\[
                          z_v|_{I_v\cap G_{F_1}}=0 .
\]
Passing to a further finite real cyclotomic layer preserves the
required nonzero root entries, the diagonal normalization and the
common \((2,0)\) frame coordinates of the target and auxiliary datum.
\end{proposition}

\begin{proof}
\proofstep{Kill the added tame inertia.}
For a new rational prime \(w\), every completion \(F_{0,v}\)
above \(w\) is \(\Q_w\). At its inertia, the old \(y\) and the
diagonal characters are trivial, and the new cocycle is a
homomorphism to an additive group killed by \(\ell\).
Write \(I_v\) for inertia and \(z\) for this restricted
additive cocycle. Since \(w\ne\ell\), wild inertia and the
prime-to-\(\ell\) factors of tame inertia have zero image. The
remaining factor is procyclic, so \(z\) factors as
\begin{equation}\label{eq:count-two-tame-quotient}
 I_v\longrightarrow I_v^{\mathrm{tame},(\ell)}\simeq\mathbb Z_\ell
  \longrightarrow\mathbb Z_\ell/\ell\mathbb Z_\ell
  \xrightarrow{\ \bar z\ }(k,+).
\end{equation} The cyclic
degree-\(\ell\) subfield of \(\Q(\mu_w)\) realizes that
quotient, since \(w\equiv1\pmod\ell\) and \(\Q(\mu_w)\) is
totally ramified at \(w\). Adjoin these subfields for all
new \(w\) to obtain \(F_1\). This extension is real and abelian
over \(\Q\); it is unramified over \(F_0\) at \(\ell\) and the
quotient of absolute groups over \(F_0\) has exponent dividing
\(\ell\).

By \eqref{eq:count-two-tame-quotient} every new inertia
restriction vanishes on \(G_{F_1}\). For \(F'=F_1\) take \(a=1\);
after adjoining a finite layer of the real cyclotomic
\(\mathbb Z_\ell\)-extension choose \(a\) so that the exponent
of \(\operatorname{Gal}(F'/F_0)\) divides \(\ell^a\). Then
\[
 \tau^{\ell^a}\in G_{F'},\qquad
 d_i(\tau^{\ell^a})-d_j(\tau^{\ell^a})
       =(d_i(\tau)-d_j(\tau))^{\ell^a}\ne0\quad(i\ne j).
\]
The restricted diagonal characters therefore still separate the labels.

\proofstep{Preserve the nonzero root entries.}
For each normalized cocycle \(u\) the normalizing element satisfies
\(u(\tau)=0\) and \(\theta(\tau)\ne1\). The quotient by the new absolute
group has exponent dividing a power of \(\ell\), so
Lemma~\ref{lem:root-cocycle-ell-power} preserves a nonzero value on
that group.
The same two-by-two multiplication applies to the retained
entries whenever the indicated block support is used.
Thus the restricted cocycle can be nonzero even if the original inertia element lies outside the subgroup.

\proofstep{Transport the local flags.}
At a place of a deeper layer, take a conjugator \(g\) over
the earlier field carrying its marked predecessor to the
place in question. Transport the flag by \(r(g)N_v\) and
take its lower unipotent factor. For \(A=r(g)N_v\) the successive determinants in the flag calculation are
\[
 \Delta_0(A)=1,\qquad \Delta_m(A)=\det(A_{ij})_{0\le i,j<m},
\]
with indices starting at zero throughout.

In the support
of the first line, the first row has no off-diagonal terms and
\(A_{21}=0\). Gaussian elimination therefore gives
\[
 \Delta_{i+1}/\Delta_i=d_i(g)\quad(0\le i\le2),\qquad
 \Delta_4/\Delta_3=\det(A)/(d_0(g)d_1(g)d_2(g)).
\]
These four pivots are units. Hence the decomposition into a lower
unipotent and an upper triangular factor is unique.
The symplectic adjoint reverses products and preserves
upper triangular and lower triangular types. Applying it
to the factorization shows by uniqueness that the lower
unipotent factor is symplectic. A different conjugator
differs on the right by an element of the original
decomposition group and hence gives the same flag.
In the transported frame the \((2,0)\) entry is
\[
             \frac{r(g)_{20}+d_2(g)(N_v)_{20}}{d_0(g)} .
\]
It is identical for the target and the auxiliary datum.
The \((2,1)\) entries remain zero. The local upper
triangular matrices have the original characters pulled
back, since conjugation inside the earlier Galois group
does not change those characters.

The same Gaussian construction carries the integral
target to each marked place. The leading principal
minors are units since their reductions are the
pivots just computed. Its flag characters are the
original rational characters restricted to inertia.
The rational point kills the inertia on the residual
kernel outside its marked set. At an original exception allowed by the target, the inertia and tower conditions give
\[
 \bar\Psi=1,\qquad\kappa=1\quad\Longrightarrow\quad\psi_0=1
\]
on the relevant inertia. Hence the rational point's other inertia condition applies there as well.
On inertia above \(\ell\) in the cyclotomic kernel,
the trivial residual labels and the prime-to-\(\ell\)
order of \(\psi_0\) on the deep tower give all four
ordered inertia characters equal to one.
\label{cor:cohomology-count-two-chain-3}
\end{proof}

\paragraph{The two Levi blocks and the three lines.}
\label{par:count-two-lines}
We can now read the path from its block decompositions. At each step the endpoints agree on the diagonal blocks. The middle step changes which
pair is held fixed.
\begin{construction}
\begin{tikzpicture}[>=Stealth,node distance=7mm,every node/.style={inner sep=2pt}]
\node (target) {Restricted rational target};
\node[below=of target] (aux) {Auxiliary datum};
\node[below=of aux] (proj) {One-root projection};
\node[below=of proj] (seed) {Mixed Hilbert seed};
\draw[->] (target) -- node[right=3mm] {\(\{0,2\},\{1,3\}\)} (aux);
\draw[->] (aux) -- node[right=3mm] {\(\{0,1\},\{2,3\}\)} (proj);
\draw[->] (proj) -- node[right=3mm] {\(\{0,2\},\{1,3\}\)} (seed);
\end{tikzpicture}
\caption{Block orders for the three affine lines of Proposition~\ref{prop:count-two-three-lines}. In the indicated order the endpoint differences belong to the square-zero off-block ideal.}\label{fig:analytic-three-lines}
\end{construction}

\begin{proposition}[Three lines joining the rational target to the seed]\label{prop:count-two-three-lines}
In the even-order case, after the restriction of
Proposition~\ref{prop:count-two-tame-return} and at a deep layer
where the mixed seed is available, the rational target and that
seed are joined by affine interpolation from target to auxiliary datum,
then from the auxiliary datum to its one-root projection, and finally
from that projection to the seed. Write \(D_0,D_1,D_2,D_3\) for the four marked data in that order.
There are families \(\mathcal D_i\) over \(k[t]\), affine in the matrix
and frame coordinates, with
\[
           \mathcal D_i(0)=D_i,\qquad \mathcal D_i(1)=D_{i+1}
                         \qquad(0\le i<3).
\]
They preserve the marked flags, similitude,
diagonal normalization and the prescribed inertia characters.
\end{proposition}

\begin{proof}
\proofstep{Project to the one-root datum.}
In the even-order case the character \(B\) is trivial on the
decomposition groups above \(\ell\). On the deeper tower,
project the auxiliary global matrices to their diagonal
blocks in the order \(\{0,1\},\{2,3\}\). Do the same
to the local frames. These projections respect products
on the block lower triangular algebra; the symplectic
form pairs the two blocks, so the projected matrices and
frames still satisfy the same similitude and flag
identities. The projected residual representation has
only the off-diagonal entries of degree \(-w_\alpha\),
including the nonzero entry in position \((1,0)\).
In its projected frame the local triangular matrix is
diagonal, because the representation and frame are
lower triangular as well.

At each marked place, fix the frame's \((2,0)\) entry using the seed value \(e_v\):
\[
 (N_v)_{20}=e_v .
\]
Right-multiply the projected frame by \(U(0,e_v,0)\).
On the marked decomposition group the first and third diagonal characters agree. Cancelling this common character in the paired-product identity gives
\[
 d_0=d_2,\qquad d_0d_3=d_1d_2\quad\Longrightarrow\quad d_1=d_3 .
\] Thus the frame multiplier commutes with the
diagonal local matrix. It preserves the \((1,0)\) entry
and inserts \(e_v\) in position \((2,0)\).

\proofstep{Construct the three lines.}
The first line joins the restricted target to the
auxiliary datum. The second joins the auxiliary datum
to this one-root projection with the modified frame.
The third joins the projection to the reduced seed.
For the first and third lines, order the blocks as
\(\{0,2\},\{1,3\}\). Both endpoints have the same
diagonal blocks, by the common \((2,0)\) entry, the
paired \((3,1)\) entry determined by the alternating
form, and the fixed diagonal characters. In the third line the two endpoints have identical readings in the global and the frame coordinates at position \((2,0)\):
\[
 (\text{global entry},\text{marked frame entry})=(0,e_v).
\] The second line uses
\(\{0,1\},\{2,3\}\), with the same diagonal blocks by
the common \((1,0)\) entry and its paired \((3,2)\)
entry.

In each order, all differences are in the
square-zero ideal of strictly off-block matrices.
The adjoint preserves this ideal, since the form
pairs the two ordered blocks. To verify the interpolation, write any interpolated matrix
as \(M_0+t\Delta M\), with \(\Delta M\) in the square-zero ideal.
For two such matrices multiplication gives
\[
 (M_0+t\Delta M)(N_0+t\Delta N)
 =M_0N_0+t(M_0\Delta N+\Delta M N_0)
 =(1-t)M_0N_0+tM_1N_1 .
\]
The adjoint is linear and preserves the ideal, so the same
calculation applies to both the product identities and the
similitude identity. Applying the equality to the individual basis columns \(e_i\) reads the marked flag along each line:
\[
 R(s)(N_v e_i)=N_vD_v(s)e_i .
\] The
interpolants are defined over \(k[t]\) and preserve
the normalization and all marked characters.

The original target vertex uses the two cases of the
lattice and restriction calculation above. The auxiliary
vertex retains the nonzero \(-w_\Psi\) if present,
and otherwise has the same nonzero marked frame entry
as the target and the new \(-w_\alpha\) entry. The
third vertex uses the nonzero \(-w_\alpha\) entry and
the seed's nonzero frame entry. For the seed, let \(S_{\ne0}\) be the set of weights of nonzero entries in its prescribed support. The two allowed alternatives are
\[
 -w_\alpha\in S_{\ne0}\qquad\text{or}\qquad \{w_\beta,-w_\vartheta\}\subseteq S_{\ne0}.
\]
In the second alternative the \(w_\alpha\)-character must also be nontrivial on the marked local group.

These choices put the auxiliary fields before the
bounds on the deeper cyclotomic tower. The finite
set of ramified auxiliary primes and \(F_1\) are fixed
using the target over \(F_0\). The cocycle bound \(C\)
is taken at this tame base \(F_1\). Choose the deeper layer only after obtaining \(C\) over \(F_1\). If \(d_{\rm layer}\) denotes its degree, impose both the depth condition and
\[
 d_{\rm layer}\ge30(C+1)^2
\]
before constructing the mixed seed on that layer. The cup lift, the ramification removal and the
transport of the target are independent of this later
seed. Its frame entry is used only in the last two lines.
\label{cor:cohomology-count-two-chain-4}
\end{proof}
